
\section{\trSymplectic\trManifold{}的定义}
\label{sec:symplectic-manifolds-definition}

我们来试着理解一下，\trSymplectic\emph{\trManifold}的定义可能会是什么样子。我们首先要求每个\trTangentSpace{}都配有上面定义的向量结构：\trManifold{} $W$ 将配有一个\trDifferential{} 2-\trForm{} $\omega$，也就是说，对每个 $x$，在 $T_xW$ 上有一个\trAlternating\trBilinearForm{} $\omega_x$。此外，我们还希望每个 $\omega_x$ 都是\trNondegenerate{}的。若 $W$ 的维数（必为偶数）是 $2n$，那么我们可以把这个条件写成
\[
\omega_x^{\wedge n}=\wedge^n\omega_x\neq 0\in\bigwedge\nolimits^{2n}T_x^\star W
\]
对每个 $x$ 都成立。换句话说，我们要求 2n-\trForm{} $\omega^{\wedge n}$ 是 $W$ 上的一个\textbf{\trVolume\trForm}（volume form）。注意，特别地，\trManifold{} $W$ 必须具有一个\trVolume\trForm，从而是\trOrientable{}的。

事实上，这个定义是不够的：我们还希望 $\omega$ 使我们得以在 $W$ 上进行的\trDifferential{}学，在局部上与我们利用那个“常值”形式在 $\RR^n\times\RR^n$ 上已经有的\trDifferential{}学完全一样。于是，让我们把这个形式写成一个\trDifferential\trForm：我们把 $\RR^n\times\RR^n$ 的向量记作 $X=(p,q)$、$X'=(p',q')$，于是便有
\[
\omega(X,X')=p\cdot q'-p'\cdot q\text{。}
\]
换句话说，
\[
\omega=\sum_{i=1}^{n}dp_i\wedge dq_i=d\Bigl(\sum_{i=1}^{n}p_i\,dq_i\Bigr)\text{。}
\]
特别地，$\omega$ 是一个\trExact\trForm。我们本可以要求\textbf{\trSymplecticForm}（symplectic form）是\trExact{}的……不巧的是，这会使得我们无法使用\trCompact\trManifold。事实上，我们有下述结果。

\begin{proposition}\label{thm:no-exact-nondegenerate-forms}
在\trCompact\trManifold{}上，不存在任何既\trNondegenerate{}又\trExact{}的 2-\trForm。\footnote{译注：原书陈述照此排印；按字面尚须默认两事：其一，$W$ 具正维数——0 维\trCompact\trManifold{}上零形式即“既\trNondegenerate{}又\trExact{}”，命题失效；其二，$W$ 无边界——带边反例即闭\trDisk{}承袭 $\RR^2$ 的 $\omega=\sum dp_i\wedge dq_i=d\bigl(\sum p_i\,dq_i\bigr)$。——译者注}
\end{proposition}

\begin{proof}
我们已经说过，$\omega$ 是\trNondegenerate{}的当且仅当 $\omega^{\wedge n}$ 是一个\trVolume\trForm。由于 $W$ 是\trCompact{}的，我们知道诸\trVolume\trForm{}都不是\trExact{}的。\footnote{例如可在 [45] 或 [74] 第一卷第八章中看到\trRham\trCohomology{}的基础知识。}现在，若 $\omega$ 是\trExact{}的，则我们就会有 $\omega=d\alpha$，但那样一来，\trVolume\trForm{} $\omega^{\wedge n}$ 就会是\trExact{}的了：$\omega^{\wedge n}=d\bigl(\alpha\wedge\omega^{\wedge(n-1)}\bigr)$。
\end{proof}

事实上，如果我们想要像在 $\RR^n\times\RR^n$ 中那样进行\emph{局部}计算，那么我们所需要的就是一个“局部\trExact{}性”条件……也就是说，一个闭性条件。这就把我们引向那个正确的定义。

\begin{definition}\label{thm:symplectic-manifold-definition}\BookIndexDefinition{T001}
一个\textbf{\trSymplectic\trManifold}（symplectic manifold）就是一个对 $(W,\omega)$，其中 $W$ 是一个\trManifold，而 $\omega$ 是一个\trNondegenerate{}的\trClosed{} 2-\trForm。这样的形式 $\omega$ 便称为\trSymplecticForm。
\end{definition}

从一个\trSymplectic\trManifold{} $(W_1,\omega_1)$ 到另一个 $(W_2,\omega_2)$ 的\trDiffeomorphism{} $\varphi$ 称为辛的，若它满足 $\varphi^\star\omega_2=\omega_1$。我们有时也说 $\varphi$ 是一个\textbf{\trSymplectomorphism}（symplectomorphism）。

当然，往后我们会常常滥用记号，在无混淆之虞时不提形式 $\omega$ 而把 $W$ 称作一个\trSymplectic\trManifold。
